% Copyright (c) 2026 by C. Klukas.
% Licensed under the MIT License. See LICENSE for details.

The preceding sections describe what \paglets{}/cpp does. This section
says what it costs, what it does not do, and where its guarantees end.
Every limit below is stated as the code enforces it today.

\subsection{The trust model}\label{sec:discussion-trust}

A mesh has four kinds of keys: admins, owners, signers and hosts
(Section~\ref{sec:mesh}). Admins decide who belongs to the mesh and what
its policy is; owners launch \paglet{}s; signers vouch for modules; hosts
run \paglet{}s. Every decision is a signed ledger record, and the mesh
trusts the ledger, not any host. A record that needs a quorum carries the
signatures of as many admins as the quorum of its epoch demands: the
genesis record sets the two numbers, admin-set records change them, and
the separate one applies to changes of the admin set itself. A stolen
admin key is removed by the remaining admins; removals win in every
conflict, and a removed admin's later records have no effect. A mesh that
must survive a stolen admin key therefore needs an admin-set quorum of at
least two; with the default of one, a single admin can remove all others.

The sandbox protects the host from the \paglet{}; nothing protects the
\paglet{} from its host. A roaming \paglet{} reaches only the imports of
its trust class, the capabilities it holds and the quotas of
Table~\ref{tab:model-limits}; by default it runs in a worker process that
the operating system confines (\code{paglets-host} falls back to its own
process when \code{paglets-worker} is not found next to it, and
\code{-{}-in-process} and \code{-{}-no-sandbox} switch the confinement off
for debugging). The confinement differs by platform: on Linux a seccomp
filter refuses the system calls that open files, create sockets or start
programs (a deny list; an allow list is planned); on macOS a sandbox
profile denies forking, executing, networking, writing files and reading
file contents; on Windows the worker runs in a job object with a
restricted token at low integrity, without processes or writes to the
user's files, but with network access, which is not restricted yet
(\code{cpp/host/src/sandbox.cpp}). A test, \code{paglets-worker
-{}-check-sandbox}, verifies on each platform that the sandbox refuses
what it must.

A host that runs a \paglet{} holds its memory image in plain form, in the
worker and, on a host with a state directory, on disk; a compromised
enrolled host can read and change the state of every \paglet{} it runs,
speak in their names (hosts stamp the sender record of every message), and
keep the \paglet{}s it receives. It cannot issue admin records
(enrollments, policy rules, module trust, admin-set changes), because it
has no admin key; the records a host key signs are enrollment requests,
grants under an allow rule, grant releases and audit entries. It cannot
read what passes through it for other hosts, because channels are
encrypted end to end. The defenses are administrative: admins revoke the
host's enrollment, owners send \paglet{}s only to hosts with a given label
(a transfer ticket names \code{label:<label>}), and admins bind marked data
to hosts by residency rules (Section~\ref{sec:mesh}). Images are not
encrypted at rest. The mesh is, as in the Python edition
\cite{klukas2026pagletsNote}, a set of hosts one organization trusts; what
is new is that one owner's \paglet{} cannot harm another's or the host,
and that the host is told what a \paglet{} may do. By default any enrolled
owner may launch any module; \code{module-policy trusted} makes roaming
modules need module trust as well, a \code{module-trust} record that names
their signer or their hash. Neither the cryptographic protocol (Noise XX
channels, record and passport signatures) nor the sandboxes have been
reviewed outside the project.

\subsection{Images, the memory limit and growth}\label{sec:discussion-memory}

A \paglet{}'s linear memory is at most 256 pages, 16\,MB
(\code{wasm::Limits::max\_memory\_pages}). The limit is a compiled-in
default of the host, not a choice of the \paglet{}. Guests are built
without exceptions (\code{-fno-exceptions}), so there is no
\code{bad\_alloc} to catch: an allocation beyond the limit aborts, an abort
is a trap, and a trap ends the \paglet{}. Messages are at most 1\,MB and
the storage service gives 16\,MB per \paglet{}. The limits are there so
that no owner can take a host over, and they shape how a \paglet{} is
written: large data belongs in files and artifacts, which move between
hosts as capabilities of a few bytes; the \paglet{} carries the handle, not
the content. The pi example keeps only the chunks that arrive early and
writes every chunk that is next in order to a file, so that 100\,000
digits need little memory.

Images have a cost that serialization does not. WebAssembly memory grows
and never shrinks, and the allocator does not zero what it frees. A
\paglet{} that once held a 10\,MB vector keeps a 10\,MB image after the
vector is gone, because the pages are no longer zero and the image elides
only zero pages. The page manifest limits what this costs in motion: a
host that has seen the \paglet{} before receives only the pages that
changed, and a clone fanned out to a host that holds the parent's pages
sends mostly hashes, because the page cache of a host is keyed by page
hash, not by \paglet{}. Image capture hashes every non-zero page with SHA-256; in the
feasibility measurements this dominated the capture time (1\,ms for two
pages on a Mac; 6.5\,ms on a Raspberry Pi 4 that was busy with a compiler
build on all four cores at the time), and the hash is still the portable
C++ implementation in \code{cpp/host/src/sha256.cpp}.

The image is bound to its module's hash. A new version of a module cannot
take over an old image, and the runtime has no state migration between
versions: a \paglet{} that must survive a code change has to hand its
state to its successor in messages, or write it to storage. A
reflection-derived export and import of state for this case is foreseen in
the plan and not built.

What images buy is the absence of serialization code and of the bugs
that come with it: nothing is forgotten in a move, pointers inside the
heap stay valid, and the \paglet{} is any C++ object graph.
Section~\ref{sec:evaluation} gives the sizes and times of real images.

\subsection{Interpreter speed}\label{sec:discussion-speed}

\paglet{}s run in WAMR's fast interpreter. AOT and JIT are turned off in
the build (\code{PagletsWamr.cmake} in \code{cpp/cmake}): no LLVM in the
host, and no compiled code that would have to be cached per CPU
architecture. The feasibility measurements in
\code{planning/cpp-m0-results.md} put the fast interpreter at about ten
times the message throughput of the classic interpreter, a host-to-guest
call at 0.4\,\textmu{}s and a message round trip through MessagePack at
6.5\,\textmu{}s on a Mac. The live benchmark of
Section~\ref{sec:evaluation} measures about 250 million integer operations
per second on a Mac Studio and 67 million on a Raspberry Pi 4, in
\paglet{} code.

For what \paglet{}s mostly do, moving, messaging and calling services, the
interpreter is not the bottleneck: a move runs guest code only in its two
hooks, and its time goes into capturing, transferring and restoring the
image, which the benchmark test of Section~\ref{sec:evaluation-images}
measures as a whole, between two hosts in one process, and does not
split. For compute the interpreter counts: the pi example needs 313\,s for
100\,000 hexadecimal digits on three hosts. How much of this is the
interpreter was not measured: no native run of the same \code{bbp.hpp} is
recorded, and the 313\,s also contain the chunk round trips across the
mesh. A \paglet{} is the right place for the coordination of such work and
a poor place for its inner loop. The way out within the model is a system
paglet that does the heavy part outside the interpreter (the \code{ai}
gateway, for instance, hands inference to an Ollama server); AOT with a
per-architecture cache, keyed by module hash and WAMR version, is the
planned alternative and was not measured.

\subsection{What weak mobility does not solve}\label{sec:discussion-weak}

A \paglet{} moves only between handlers. There is no migration in the
middle of a computation: no stack is captured, and a handler that runs
longer than the budget of 5\,s is stopped and its \paglet{} ends. Every
long computation must therefore be sliced by its author. The pi example
converts its hexadecimal fraction into decimal digits in slices of
200\,ms, each started by a timer of zero delay, and keeps its progress in
members between the slices (Listing~\ref{lst:pi-convert}). In the listing
the loop condition \code{now\_ms() - began < 200} bounds the slice, the
members \code{limbs\_} and \code{decimal\_done\_} carry the progress, and
\code{after\_raw(0, "convert", \{\})} near the end schedules the next
slice; everything in between is the arithmetic. Between two slices the
\paglet{} can be moved, checkpointed or deactivated; inside a slice it
cannot.

% snippet: cpp/examples/pi/pi.cpp#convert [float,caption={Pi's decimal conversion, sliced into handler calls of 200\,ms},label={lst:pi-convert}]
\begin{lstlisting}[style=paglets-cpp,float,caption={Pi's decimal conversion, sliced into handler calls of 200\,ms \src{cpp/examples/pi/pi.cpp}},label={lst:pi-convert}]
// A slice of the conversion: passes of multiplying the fraction by 10^9,
// whose overflow above the point is the next nine decimal digits.
void convert() {
    if (state_ != "converting") return;
    const std::int64_t began = now_ms();
    std::string digits;
    while (decimal_done_ + static_cast<std::int64_t>(digits.size()) < decimal_total_ && now_ms() - began < 200) {
        // Only the precision the remaining digits need (and two limbs more).
        const std::int64_t remaining = decimal_total_ - decimal_done_ - static_cast<std::int64_t>(digits.size());
        const auto need = static_cast<std::size_t>(static_cast<double>(remaining) * 3.3219280949 / 32.0) + 2;
        if (limbs_.size() > need) limbs_.resize(need);
        std::uint64_t carry = 0;
        for (std::size_t i = limbs_.size(); i-- > 0;) {
            const std::uint64_t t = static_cast<std::uint64_t>(limbs_[i]) * 1'000'000'000ULL + carry;
            limbs_[i] = static_cast<std::uint32_t>(t);
            carry = t >> 32;
        }
        char nine[16];
        std::snprintf(nine, sizeof nine, "%09u", static_cast<unsigned>(carry));
        digits.append(nine, static_cast<std::size_t>(std::min<std::int64_t>(9, remaining)));
    }
    decimal_done_ += static_cast<std::int64_t>(digits.size());
    if (shown_dec_.size() < shown_digits) shown_dec_ += digits.substr(0, shown_digits - shown_dec_.size());
    if (dir_) {
        dec_file_.pending += digits;
        flush(dec_file_);
    }
    if (decimal_done_ < decimal_total_) {
        (void)paglets::after_raw(0, "convert", {});
        return;
    }
    std::vector<std::uint32_t>().swap(limbs_);
    state_ = "writing";
    maybe_done();
}
\end{lstlisting}
% /snippet

The slicing is visible in the code and is the author's responsibility; the
runtime cannot do it for a loop it does not see. Suspending a handler in
the middle (Asyncify or Wasm stack switching; WP-R1 in the plan) would
remove the slicing. It would not change what moves: a pending service
call, a timer or a reply capability lives on the host, not in linear
memory, and each follows its own rule on a move
(Section~\ref{sec:movement}) whether the handler was suspended or not.

\subsection{Platforms and editions}\label{sec:discussion-platforms}

\paglets{}/cpp is built and tested in CI on Linux x86-64 and arm64, on
macOS arm64, and under AddressSanitizer, UndefinedBehaviorSanitizer and
ThreadSanitizer. Windows is experimental: the job builds with MinGW-w64
GCC in MSYS2, runs with guest modules built by the Linux job, uses
software bound checks because MinGW cannot build WAMR's guard-page
checks, and may fail without failing the workflow
(\code{continue-on-error} in \code{.github/workflows/cpp.yml}). The
worker sandbox on Windows is the weaker one described above. On the
Raspberry Pi, kernels with a 39-bit address space hold about 60 instances
per process with guard-page checks, because each instance reserves 8\,GB
of address space. Configuring with
\code{PAGLETS\_WAMR\_HW\_BOUND\_CHECK=OFF} removes the limit (the default
is on everywhere but MinGW); software bound checks made message handling
about 2.5 times slower on macOS in the feasibility measurements.

The two editions of \paglets{} do not interoperate. They share concepts
(hosts, \paglet{}s, messages, meshes, services, compute slots) and no
format: transport, authentication, framing, movement and code differ
throughout \cite{pagletsCppDocs}. A Python host and a C++ host cannot talk
to each other, and a \paglet{} of one edition cannot run on the other.

\subsection{Three systems compared}\label{sec:discussion-comparison}

Table~\ref{tab:background-editions} placed the three systems next to
each other by what moves and how it runs;
Table~\ref{tab:discussion-editions} adds what each one pays for it. The
price of \paglets{}/cpp is the one this section
described: a fixed memory limit, an interpreter, sliced handlers, one
guest language, and a trust model that needs administration. What it buys
over the Python edition is that hosts need no application code, that a
\paglet{} cannot reach past its capabilities, and that a job survives the
loss of hosts.

\begin{table}[htbp]
\centering\small
\caption{What each system pays for an agent.}
\label{tab:discussion-editions}
\begin{tabular}{p{0.16\linewidth}p{0.23\linewidth}p{0.23\linewidth}p{0.25\linewidth}}
\toprule
 & \raggedright Aglets \cite{lange1997aglets,ferrari2004manual} & \raggedright \paglets{}, Python \cite{klukas2026pagletsNote} & \raggedright \paglets{}/cpp \tabularnewline
\midrule
\raggedright State on the wire & \raggedright the Java object, serialized; classes with it, or fetched from the codebase & \raggedright a dataclass, serialized; no code & \raggedright the non-zero 64\,KB pages the destination lacks, compressed; modules by SHA-256, fetched on demand \tabularnewline
\raggedright Memory per agent & \raggedright a share of the server's JVM heap & \raggedright a Python process per active \paglet{} & \raggedright at most 16\,MB of linear memory; an image keeps every page that was ever non-zero \tabularnewline
\raggedright Execution & \raggedright the JVM & \raggedright CPython & \raggedright a Wasm interpreter, no JIT or AOT \tabularnewline
\raggedright Long computations & \raggedright a thread per aglet & \raggedright a handler in the \paglet{}'s process; its mailbox waits & \raggedright sliced handlers; a budget of 5\,s per handler \tabularnewline
\raggedright Fault tolerance & \raggedright --- & \raggedright deactivated \paglet{}s come back from disk after a host restart; active ones are lost & \raggedright restart from the last image; zombie replicas on every host that offers \code{zombies} \tabularnewline
\raggedright Sandbox & \raggedright the Java platform & \raggedright none; a \paglet{} reaches what its process reaches & \raggedright Wasm instance, import set per trust class, sandboxed worker processes \tabularnewline
\raggedright Administration & \raggedright a Tahiti server per host & \raggedright a shared API key & \raggedright keys for every party, a signed ledger with quorums, policy, module trust, residency \tabularnewline
\raggedright Tooling & \raggedright the Tahiti server & \raggedright a CLI & \raggedright CLI, web UI, interface descriptions, protocol checker \tabularnewline
\bottomrule
\end{tabular}
\end{table}

\subsection{Open problems}\label{sec:discussion-open}

\begin{itemize}
\item \emph{Verification of the zombie protocol.} The exhaustive check is
  complete for two replicas and bounded to a depth of 11 steps for three
  (Section~\ref{sec:zombies}). Every world is explored with hash
  compaction (states identified by a 128-bit hash of their encoding,
  \code{verify/explorer.hpp}), so even the two-replica results carry a
  small omission probability, and the timed random runs sample timer
  values and takeover times instead of enumerating them. An exact mode, a
  fairness check of the fault-free graph, a mutation suite that shows the
  checker catches seeded bugs, and either a full three-replica exploration
  or a cutoff argument are planned; the repository tracks these items. The
  model also leaves out the transport, the gossip and the \code{zombies}
  service itself; the replay of recorded live traces is the conformance
  check between them.
\item \emph{Mesh size.} The live mesh has three hosts. The tests run
  meshes of four hosts in one process and a gossip test of five hosts
  under random partitions. Nothing has been measured beyond that: how the
  location records, the ledger gossip and \code{mesh-info} snapshots
  behave with tens of hosts is not known.
\item \emph{Relays and NAT.} Hosts without an inbound port take part
  through relays, and the test suite exercises this with hosts that do not
  listen, including two such hosts reaching each other and a relay going
  down (\code{cpp/tests/test\_discovery.cpp}). No live run has crossed a
  NAT or a firewall; the live hosts share one LAN.
\end{itemize}

The limits of the sections above are open as well: the security review,
state migration between module versions (foreseen in the plan, not
built), AOT compilation with a per-host cache, and the research items
WP-R1 (mid-handler migration) and WP-R2 (other guest languages: the ABI
is core WebAssembly, so a second language is a matter of an SDK, not of
the host).
